% !TeX root = main.tex
\documentclass[12pt]{article}

\usepackage{tikz}

\title{An investigation on the relationship between slit distance and fringe spacing}
\author{Jimmy Zeng}

\begin{document}
\maketitle
\tableofcontents
\section{Introduction and Background}
Wave Phenomena states that wave could interfere when they overlap. When two waves are in opposite phase, it is said to interfere destructively and cancel out. However, if they are the same phase and interferes, it is said to be constructive, reaching a greater amplitude.
% Include Diagram
The debate of whether light is a particle or a wave was challenged by Thomas Young's famous Double Slit Experiment in Young's Course of Lectures on Natural Philosophy and the Mechanical Arts of 1807, where a coherent light source shinning through two tiny slits projected an interference pattern with alternation between light and dark fringes rather than two individual bright fringes in a dark background. The formula proposed by Thomas Young was
%Include Citations
\begin{equation}
    s=\frac{\lambda D}{d}
\end{equation}
in Which "S" is the distance between two adjacent bright fringes, $\lambda$ is the wavelength of the light, $D$ is the distance between the screen and slits in meters and $d$ is the distance between slits. 
The mathematical geometric derivation is thus justified: %\cite{}
% Include mathematical Derivation with diagrams
Take a look at the following graph:
%put the forking graph here
We observe that the path difference from the wave from the bottom slit and top slit is illustrated as $\Delta x$, this coincides with the equation
\begin{equation}
    d sin(\theta) = \Delta x
\end{equation}
We observe the fact that the path difference ($\Delta x$) for constructive waves is
\begin{equation}
    n\lambda
\end{equation}
where $\lambda$ is the wavelength of the light and $n$ is an integer representing the interested bright fringe.

%Use diagrams to better explain and stuff

\section{Methodology}
    

\subsection{Research Hypothesis and Method variables}

The independent variable is the distance between the slit and the screen and will include values that range from the [nearby and further apart]. The [nearby] will include 20cm, 40cm, 60cm, 80cm, and 100cm, while the further apart include 3 meters, 5 meters, 6 meters and 8 meters.
% Include diagram from chemix to elucidate

Since the fringe spacing is within the milimeter range and it is beyond my human ability to hold the clipboard still within milimeter stillness while using the other hand to hold a pencil and mark accurately on the paper live. I decided to take a high resolution photo of the projection on a graphed piece of paper as a photo is most versatile and captures exactly as seen. The graphed piece of paper is printed with milimeter accuracy markings. As a result I will measure distance between each adjacent fringe (in milimeters) with software solutions and excuse my skill issue.

Our guiding question is \textbf{What is the relationship with the distance between screen and slit (from the range of 20cm to 8 m) and the spacing between adjacent bright fringes (in mm) in a double slit system.}


    


\end{document}


